% ch07-gof — goodness of fit, model comparison, and calibration discipline

\section{The evaluation constitution}\label{ch07:sec:constitution}

Every preceding chapter builds an estimator; this chapter fixes the court in
which all of them are tried. Its deliverables are not models but rules: a
definition of ``calibrated'' for each of the four targets of
\cref{ch:question}, the test batteries that operationalize those
definitions, the comparison ladder that decides which model family is
promoted, and the reporting discipline that keeps every number attached to
the caveats it needs. We call this the \emph{evaluation constitution}
because it is meant to be hard to amend: metrics chosen after seeing
results are not metrics, and the repository's synthetic rehearsals
(\cref{app:results}) demonstrated repeatedly that the difference between a
defensible claim and an embarrassing one was almost never the model --- it
was whether the evaluation respected the risk set, the temporal split, and
the identifiability limits of the parameter being quoted.

The word ``calibrated'' means something different for each target, and
conflating the meanings is the failure mode
\cref{ch01:sec:targets} warns against.

\begin{definition}[Calibration, per target]\label{def:ch07:calibrated}
Fix an evaluation era and the point-in-time discipline of
\cref{prot:ch02:pit}.
\begin{enumerate}[label=(\roman*)]
\item \emph{T2 (timing).} An intensity or count model is calibrated if it
survives the time-rescaling battery of \cref{prot:ch07:rescaling} at the
finest resolution its event times honestly support, \emph{and} the
count-calibration battery of \cref{prot:ch07:pit} (PIT uniformity,
dispersion near $1$, reliability per sector-week) at weekly resolution.
Both are required: rescaling sees timing structure that counts hide, and
PIT sees marginal-distribution failures that rescaling averages out.
\item \emph{T3 (horizon events).} A horizon classifier is calibrated if its
predicted probabilities match conditional frequencies:
$\Prob(\text{event within } \horizon \st \hat p = p) = p$ across the
support of $\hat p$, assessed by reliability diagram, calibration
intercept and slope. Discrimination (PR-AUC, lift at $K$) is necessary
but never sufficient; recalibration maps
\citep{platt1999,zadroznyelkan2002,niculescumizil2005} are legal but are
model components, fitted on the validation era only.
\item \emph{T1 (identity).} A ranker is judged by ranking metrics
(recall at $K$, MRR, computed on the declared risk set) \emph{and} by
held-out conditional likelihood (nats per deal of the mark factor). Each
alone is gameable: recall at $K$ is blind to the probabilities attached to
the ranks it gets right, and conditional NLL can improve while the tail
ordering that sourcing consumes inverts. Both must be reported; neither
may be traded silently for the other.
\item \emph{T4 (marks).} Mark models are scored by proper scoring rules per
mark (\cref{ch07:sec:counts}), with the same temporal split.
\end{enumerate}
\end{definition}

\begin{decision}{D7.1 --- the evaluation constitution}
The definitions in \cref{def:ch07:calibrated}, the protocols
\cref{prot:ch07:rescaling,prot:ch07:pit,prot:ch07:recovery,prot:ch07:temporal},
and the ladder of \cref{ch07:sec:comparison} are the only admissible
evaluation for models in this repository. A new metric enters the
constitution only together with a new consuming action, per decision D1.1
of \cref{ch:question}; a metric may never be added because a model
happens to do well on it. Rationale: post hoc metric selection is the
cheapest way to fool ourselves, and the rehearsal campaigns showed that
every family looks good on \emph{some} score. Reversal trigger: a change
in the decision layer (\cref{ch:decision}) that adds or retires a
consuming action.
\end{decision}

\section{Continuous-time goodness of fit}\label{ch07:sec:continuous}

\subsection{Time rescaling in practice}\label{ch07:sec:rescaling}

The instrument for event-time fit is the time-rescaling theorem, proved in
\cref{ch:foundations} from Watanabe's characterization and restated here
in the form the diagnostics use \citep{brown2002,ogata1988,daley2003}.

\begin{theorem}[Time rescaling, evaluation form]\label{thm:ch07:rescale}
Let $\counting$ be a point process on $[0,\Tend]$ with event times
$t_1 < t_2 < \cdots < t_n$, adapted to $\{\Ft\}$, admitting a strictly
positive conditional intensity $\lam(t)$ whose compensator
$\Lam(t)=\int_0^t \lam(u)\dd u$ is a.s.\ finite and continuous. Then the
rescaled times $\tau_m = \Lam(t_m)$ form a unit-rate Poisson process on
$[0,\Lam(\Tend)]$; equivalently, with $\tau_0=0$, the gaps
\begin{equation}\label{eq:ch07:gaps}
\xi_m \;=\; \tau_m - \tau_{m-1} \;\sim\; \mathrm{Exp}(1)
\quad\text{i.i.d.},
\qquad
u_m \;=\; 1 - e^{-\xi_m} \;\sim\; U(0,1)
\quad\text{i.i.d.}
\end{equation}
\end{theorem}

\begin{proof}
Since $\lam>0$, $\Lam$ is continuous and strictly increasing, so
$s\mapsto\Lam^{-1}(s)$ is a well-defined family of $\{\Ft\}$-stopping
times. Set $\tilde\counting(s)=\counting(\Lam^{-1}(s))$ and
$\tilde\F_s = \F_{\Lam^{-1}(s)}$. Because
$\counting(t)-\Lam(t)$ is a local martingale
(\cref{ch:data}, \cref{eq:ch02:martingale}), optional stopping at
$\Lam^{-1}(s)$ shows $\tilde\counting(s)-s$ is a local
$\{\tilde\F_s\}$-martingale: $\tilde\counting$ is a counting process with
deterministic continuous compensator $s$, hence a unit-rate Poisson
process by Watanabe's characterization (\cref{ch:foundations}). The gap
and uniform statements are the standard exponential-spacings and
probability-integral facts. The proof fails exactly where its hypotheses
do: if $\lam$ vanishes on at-risk gaps, rescale over at-risk time only;
and if the $t_m$ are jittered (defect (C4)), the theorem still holds for
the \emph{true} times but the computed $\tau_m$ use corrupted ones --- the
test then interrogates model and jitter jointly
(\cref{ch07:sec:reporting}).
\end{proof}

In practice $\Lam$ is the \emph{fitted} compensator $\hat\Lam$, and three
details separate a usable diagnostic from a misleading one: the open final
interval, the null distribution under estimated parameters, and
localization.

\begin{protocol}[Time-rescaling battery]\label{prot:ch07:rescaling}
\leavevmode
\begin{enumerate}[label=(\alph*)]
\item \emph{Residuals.} Compute $\xi_m$ from \cref{eq:ch07:gaps} with the
fitted compensator, integrating over at-risk time only, per sector and
pooled (\modrepo{mbpp/gof.py}, \code{time\_rescaling\_residuals}).
\item \emph{The last interval is censored.} The final gap
$\hat\Lam(\Tend)-\tau_n$ is a right-censored $\mathrm{Exp}(1)$ draw by
memorylessness. It is either excluded from the KS sample or included as a
censored observation in a likelihood-based check; it is \emph{never}
treated as a complete gap. Treating it as complete injects one draw that
is stochastically too small per unit --- across many sectors this
manufactures a spurious excess of short gaps, i.e.\ phantom clustering.
\item \emph{KS and QQ.} Test the complete gaps against $\mathrm{Exp}(1)$
(KS, \code{ks\_test\_exp1}) and plot QQ coordinates (\code{qq\_exp1})
with envelope bands. Report the statistic per sector and pooled.
\item \emph{Localization.} A rejection is the beginning, not the end:
identify \emph{which} sectors, which calendar windows, and which
covariate strata carry the misfit, using the stratified residuals of
\cref{eq:ch07:mres} --- this is Ogata's residual-analysis program
\citep{ogata1988}.
\item \emph{Defect exposure.} Under (C4) this battery is run only at the
resolution the dates honestly support (daily-resolved deals in the
event-time track; never sub-week on jittered dates); under (C3) the
fitted $\hat\lam$ absorbs dead exposure and the test inherits the
optimism of \cref{warn:ch07:deadgof}.
\end{enumerate}
\end{protocol}

\begin{warning}[Estimated parameters make standard bands optimistic for the model]\label{warn:ch07:estimated}
The KS null distribution assumes a fixed hypothesized law. When the
compensator is fitted on the same data, $\hat\Lam$ adapts to the sample,
the KS statistic is stochastically smaller than under a fixed truth, and
standard $p$-values are conservative: a model can \emph{pass} while
wrong. The honest null is simulation-based: simulate from the fitted
model, refit, recompute the statistic, and compare the observed value to
this refit-in-loop distribution. The envelope machinery of
\cref{prot:ch07:envelope} provides the harness; held-out-era residuals
(fit on train, rescale on test) are the cheap partial fix and are the
repository default.
\end{warning}

\subsection{Martingale residuals over strata}\label{ch07:sec:mres}

Time rescaling compresses the whole fit into one distributional statement;
martingale residuals decompose it. For any predictable stratum $g$ --- a
set of (unit, time) pairs selected by sector, calendar window, or
covariate values known at $t^-$ --- define
\begin{equation}\label{eq:ch07:mres}
R(g) \;=\; \sum_{m}\ind{(e_m,t_m)\in g}
 \;-\; \sum_{e}\int_0^{\Tend} \ind{(e,u)\in g}\,\atrisk_e(u)\,
 \hat\lam_e(u)\dd u,
\qquad
Z(g) \;=\; \frac{R(g)}{\sqrt{\hat\Lam(g)}},
\end{equation}
where $\hat\Lam(g)$ is the integral term. Under a correctly specified
model $R(g)$ is the terminal value of a mean-zero martingale with
predictable variation $\Lam(g)$, so $Z(g)$ is approximately standard
normal when $\hat\Lam(g)$ is large \citep{abgk1993}. The design use is
\emph{systematic lack-of-fit localization}: compute $Z(g)$ over a fixed
grid of strata --- sector $\times$ calendar quarter, staleness-age bands
of $\staleage_i(t)$, macro-covariate terciles --- and read the sign
pattern. A block of positive $Z$ across all sectors in one era is a
missing calendar covariate; positive $Z$ concentrated at high staleness
age is the (C5) signature that \cref{ch:stale} predicts; one wild cell is
usually data hygiene, not model failure. Isolated $|Z|>2$ values among
dozens of strata are expected; structure is what indicts the model.

\subsection{Simulation-based envelope tests}\label{ch07:sec:envelope}

Likelihood-based diagnostics ask whether the data look probable under the
model; envelope tests ask whether the model's own sample paths look like
the data. They are the only GOF instrument that sees \emph{functionals we
did not fit}.

\begin{protocol}[Envelope battery]\label{prot:ch07:envelope}
\leavevmode
\begin{enumerate}[label=(\alph*)]
\item Simulate $B$ (default $999$) replicates of the evaluation era from
the fitted model using the simulators of \cref{ch:hawkes}: Ogata thinning
\citep{ogata1981}, the cluster representation \citep{hawkesoakes1974}, or
exact exponential-kernel simulation \citep{dassioszhao2013}; pass each
replicate through the same observation layer (aggregation, at-risk masks)
as the real data.
\item Compute, per replicate and for the data, a fixed pre-registered
statistic vector: weekly-count autocorrelations at lags $1$--$8$, the
dispersion index of \cref{prop:ch07:dispersion}, the coefficient of
variation of inter-event gaps, the maximum weekly count, and the
sector-level count cross-correlations.
\item Reject a statistic at (pointwise) level $\alpha \approx 2/(B+1)$ if
the observed value falls outside the central simulated band; report all
statistics together and treat the multiplicity honestly --- one excursion
among a dozen statistics at $B=999$ is weak evidence.
\item \emph{Defect exposure.} Simulating through the observation layer is
mandatory: envelopes from the latent (uncorrupted) model test the wrong
null. (C1)/(C4) are reproduced by construction; (C3) cannot be reproduced
(we do not know who is dead) and the envelope inherits the same
contaminated risk set as the fit.
\end{enumerate}
\end{protocol}

The double-use caveat of \cref{warn:ch07:estimated} applies verbatim:
parameters estimated from the data pull the envelope toward the data, so
envelope non-rejection is weak; envelope \emph{rejection} is strong and
localizes what the model cannot reproduce (typically: the real data's
overdispersion or its cross-sector comovement).

\section{Count-level goodness of fit}\label{ch07:sec:counts}

At weekly resolution the model emits a predictive distribution
$F_{s,w}$ for each sector-week count $Y_{s,w}$, and the calibration
question becomes: are the observed counts distributed as the model says?
The continuous-data answer --- the probability integral transform
$F(Y)\sim U(0,1)$ --- fails for discrete data: $F(Y)$ is a discrete random
variable, stochastically larger than uniform, and naive PIT histograms
read miscalibration where there is none, most severely in sparse cells
where $\Prob(Y=0)$ is large. Two standard repairs exist
\citep{smith1985,brockwell2007,czado2009}.

\begin{proposition}[Randomized PIT]\label{prop:ch07:rpit}
Let $Y$ have cumulative distribution $F$ on $\N_0$ with
$p_y=F(y)-F(y-1)$, and let $V\sim U(0,1)$ independent of $Y$. Then
\begin{equation}\label{eq:ch07:rpit}
U \;=\; F(Y-1) + V\,p_Y
\end{equation}
satisfies $U \sim U(0,1)$.
\end{proposition}

\begin{proof}
Fix $u\in(0,1)$ and let $y^{*}$ be the unique value with
$F(y^{*}-1) < u \le F(y^{*})$. Conditioning on $Y=y$:\ if $F(y)\le
u$ (i.e.\ $y<y^{*}$) then $U\le u$ with probability $1$; if $y>y^{*}$
then $F(y-1)\ge u$ and $U\le u$ has probability $0$; if $y=y^{*}$ then
$\Prob(U\le u \st Y=y^{*}) = (u - F(y^{*}-1))/p_{y^{*}}$. Summing,
\[
\Prob(U\le u)
 = \sum_{y<y^{*}} p_y + p_{y^{*}}\cdot\frac{u-F(y^{*}-1)}{p_{y^{*}}}
 = F(y^{*}-1) + u - F(y^{*}-1) = u. \qedhere
\]
\end{proof}

Under a correctly calibrated sequence of predictive distributions the
$U_{s,w}$ are i.i.d.\ uniform and the whole continuous toolkit (KS,
histograms, drift plots) applies. The price is randomization: two
analysts get two dashboards. The non-randomized alternative of
\citet{czado2009} replaces each draw by its conditional expectation: with
$a=F_{s,w}(y_{s,w}-1)$, $b=F_{s,w}(y_{s,w})$, define the conditional PIT
$c(u;a,b) = \min\{\max\{(u-a)/(b-a),0\},1\}$ and the aggregate
$\bar F(u) = \frac{1}{n}\sum_{s,w} c(u;a_{s,w},b_{s,w})$; calibration
implies $\E\,\bar F(u)=u$, and the histogram of increments of $\bar F$ is
the deterministic analogue of the PIT histogram. Repository rule: the
non-randomized version drives dashboards and regressions (reproducible),
the randomized version drives formal tests (exact uniformity).

\begin{proposition}[Dispersion index and its sampling distribution]\label{prop:ch07:dispersion}
Let $Y_i \sim \mathrm{Poisson}(\mub_i)$ independently, $i=1,\dots,n$, and
\begin{equation}\label{eq:ch07:dispersion}
\hat\phi \;=\; \frac{1}{n-p}\sum_{i=1}^{n}
 \frac{(Y_i-\hat\mub_i)^2}{\hat\mub_i},
\end{equation}
with $p$ fitted parameters. With known means ($p=0$, $\hat\mub_i=\mub_i$):
$\E\,\hat\phi = 1$ and
$\Var\big(\hat\phi\big) = n^{-2}\sum_i \big(2 + 1/\mub_i\big)$, so that if
$n^{-1}\sum_i 1/\mub_i \to \bar m < \infty$,
\[
\sqrt{n}\,\big(\hat\phi - 1\big) \;\xrightarrow{d}\;
\mathcal{N}\!\big(0,\; 2 + \bar m\big).
\]
Estimated means reduce the degrees of freedom, whence the divisor $n-p$.
\end{proposition}

\begin{proof}
For $Y\sim\mathrm{Poisson}(\mub)$ the central moments are
$m_2=\mub$, $m_4=\mub+3\mub^2$. Hence
$\E[(Y-\mub)^2/\mub]=1$ and
$\Var[(Y-\mub)^2/\mub] = (m_4-m_2^2)/\mub^2
 = (\mub+2\mub^2)/\mub^2 = 2 + 1/\mub$. Independence gives the variance
sum; the Lindeberg CLT (bounded $1/\mub_i$ suffices) gives normality.
The proof shows what practitioners forget: the null variance is
\emph{not} $2/n$ --- the $1/\mub_i$ term dominates in sparse cells, and
sector-weeks with $\mub_i \lesssim 1$ (common in this data) widen the
honest acceptance band. Judging $\hat\phi$ against
$1 \pm 2\sqrt{2/(n-p)}$ in a sparse panel manufactures false
overdispersion alarms.
\end{proof}

Overdispersion ($\hat\phi>1$) is the escalation trigger of the model
ladder --- it says history or heterogeneity is missing --- but it is
\emph{never} evidence for excitation specifically, because frailty and
contagion produce it identically (\cref{ch:hawkes},
\cref{ch07:sec:identifiability}).

\begin{protocol}[Count-calibration battery]\label{prot:ch07:pit}
For every count model, on the evaluation era: (a) non-randomized PIT
diagram pooled and per sector; (b) randomized-PIT KS test, seeds logged;
(c) dispersion \cref{eq:ch07:dispersion} with the honest band of
\cref{prop:ch07:dispersion}, reported \emph{alongside every count NLL};
(d) reliability per sector-week: bin cells by predicted mean decile,
compare observed to predicted bin means with Poisson error bars;
(e) PIT drift over calendar time, plotted --- drift is the cheapest
detector of vendor-coverage change and of (C5) staleness confounding.
Defect exposure: complete sector counts make this battery robust to (C6)
at sector level; (C3) does not bite sector counts (sectors do not die);
(C4) is absorbed by construction --- weekly counts are the honest
resolution.
\end{protocol}

\paragraph{Scoring rules for counts.}
A scoring rule $S(F,y)$ assigns a loss to forecasting $F$ when $y$
materializes; it is \emph{proper} if $\E_{Y\sim G}\,S(G,Y) \le
\E_{Y\sim G}\,S(F,Y)$ for all $F$, and \emph{strictly proper} if equality
forces $F=G$ \citep{gneitingraftery2007}. Propriety is not decoration: an
improper score rewards hedged or sharpened forecasts and corrupts every
comparison built on it. We use two:
\begin{equation}\label{eq:ch07:scores}
\mathrm{logS}(F,y) \;=\; -\log f(y),
\qquad
\mathrm{RPS}(F,y) \;=\; \sum_{k=0}^{\infty}\big(F(k) - \ind{y\le k}\big)^2 .
\end{equation}

\begin{proposition}[Strict propriety of the log score]\label{prop:ch07:logscore}
On the class of probability mass functions with common support,
$\mathrm{logS}$ is strictly proper.
\end{proposition}

\begin{proof}
Let $g$ be the true pmf and $f$ any competitor with the same support.
Then
\[
\E_g\big[-\log f(Y)\big] - \E_g\big[-\log g(Y)\big]
 \;=\; \sum_y g(y)\log\frac{g(y)}{f(y)}
 \;=\; \KL{g}{f} \;\ge\; 0,
\]
with equality iff $f=g$, by Jensen's inequality applied to the strictly
concave logarithm ($\sum g\log(f/g) \le \log\sum f = 0$, strict unless
$f/g$ is constant). Hence the expected score is uniquely minimized at the
truth. If supports differ the score can be infinite --- which is a
feature: a model that assigns zero mass to an observed count is
infinitely wrong, and the repository treats $f(y)=0$ on observed data as
a bug, not an outlier.
\end{proof}

The ranked probability score is proper on distributions with finite mean
\citep{gneitingraftery2007,czado2009} but not local: it rewards mass
\emph{near} the outcome. The two scores are complementary --- logS is the
likelihood and dominates the ladder (\cref{ch07:sec:comparison}); RPS is
distance-sensitive and robust to a single mis-tailed cell --- and both are
reported, logS as primary, RPS as the sanity companion.

\section{Model comparison}\label{ch07:sec:comparison}

\subsection{The held-out NLL ladder}\label{ch07:sec:ladder}

The primary comparison metric of the repository is held-out negative log
likelihood, in two units: \emph{nats per event} for event-time models
($-\loglik/n$ over test-era events) and \emph{nats per sector-week cell}
for count models ($-\frac{1}{|\text{cells}|}\sum_{s,w}\log
\hat f_{s,w}(y_{s,w})$). Every model is reported as a difference against
three named baselines fitted on the same train era: \emph{random}
(homogeneous Poisson, one global rate), \emph{historical-mean}
(per-sector constant rate), and \emph{seasonal NHPP} (the
covariate/seasonal Poisson model of \cref{ch:poisson}). A model that has
not been run against all three has not been evaluated.

\begin{warning}[Units are not exchangeable]\label{warn:ch07:units}
Nats per event and nats per cell are different currencies: they condition
on different information and normalize by different denominators. No
number in one unit may be compared to a number in the other. Comparing an
event-time model to a count model requires evaluating both at a common,
honest resolution --- in practice, aggregating the event-time model's
predictive distribution to weekly cells and paying the aggregation cost
explicitly (\cref{ch:censoring}).
\end{warning}

\begin{decision}{D7.2 --- NLL ladder with the covariates-dominate prior}
Held-out NLL against the three named baselines is the primary promotion
metric for T2 models, with \cref{prot:ch07:temporal} governing the split.
The standing prior, inherited from the rehearsal campaigns
(\cref{app:results}), is that \emph{covariates dominate and excitation is
decimal points}: moving from the seasonal baseline to the covariate NHPP
was worth $+3.5$/$+5.7$/$+1.5$ nats across the three rehearsal worlds,
while adding weekly excitation contributed $+0.04$/$-1.15$/$+0.53$.
Consequently, any claim that excitation matters must clear the seasonal
NHPP \emph{and} the covariate NHPP under rolling-origin evaluation, not
merely improve on a weak baseline. Reversal trigger: excitation
contributing $\ge 0.5$ nats per cell over the covariate NHPP on real data
in two consecutive campaigns retires the prior.
\end{decision}

AIC and BIC have a narrow legal scope here: quick screening \emph{within}
a family, on the train era, when models differ by a few regular interior
parameters. Three limits keep them off headlines. First, they are
in-sample devices, and our risk is out-of-era: a temporal regime shift is
invisible to them. Second, their penalty calculus assumes regular
asymptotics, which fail exactly at the comparison we care most about ---
excitation present versus absent --- because $\branch=0$ is a boundary
point, where even the effective parameter count is wrong (below). Third,
under misspecification neither estimates what its derivation promises.
Held-out NLL is the arbiter; information criteria are a screen.

\subsection{Testing for excitation at the boundary}\label{ch07:sec:boundary}

The natural test of ``is there any excitation?'' is the likelihood ratio
between the covariate NHPP ($\branch=0$) and the exponential-Hawkes
extension ($\branch\ge 0$). But $\branch=0$ lies on the boundary of the
parameter space --- negative branching is not in the model --- so the
classical $\chi^2_1$ asymptotics fail \citep{selfliang1987}.

\begin{theorem}[Chi-bar-square limit for one bounded parameter]\label{thm:ch07:chibar}
Let the model have parameter $(\branch,\psi)$ with $\branch\ge 0$ on the
boundary at the null, $\psi$ interior nuisance parameters, and the decay
$\decay$ \emph{fixed} (a decay bank, \cref{ch:hawkes}). Assume the LAN
regularity conditions of \citet{selfliang1987}: consistency, a
nonsingular Fisher information $\Fisher$ at the true value
$(\,0,\psi_0)$, and differentiability in quadratic mean in a neighborhood
intersected with the cone $\{\branch\ge0\}$. Then the likelihood-ratio
statistic $\lambda_n = 2\big(\sup_{\branch\ge0,\psi}\loglik -
\sup_{\psi}\loglik|_{\branch=0}\big)$ satisfies
\begin{equation}\label{eq:ch07:chibar}
\lambda_n \;\xrightarrow{d}\;
\tfrac12\,\chi^2_0 + \tfrac12\,\chi^2_1,
\qquad\text{i.e.}\qquad
\Prob(\lambda_n \le x) \to \tfrac12 + \tfrac12\,\Prob(\chi^2_1\le x),
\end{equation}
where $\chi^2_0$ is the point mass at zero.
\end{theorem}

\begin{proof}
Profile out $\psi$: under LAN, uniformly over
$\branch = O(n^{-1/2})$,
\[
\loglik(\branch) - \loglik(0)
 \;=\; \branch\,\tilde S_n \;-\; \tfrac12\,\branch^2\,\tilde\Fisher_n
 \;+\; o_p(1),
\]
where $\tilde S_n$ is the efficient score for $\branch$ (the score
residualized on the nuisance scores) and $\tilde\Fisher_n$ the efficient
information, with $Z_n \defeq \tilde S_n/\tilde\Fisher_n^{1/2}
\xrightarrow{d} \mathcal{N}(0,1)$ at the null. Maximizing the quadratic
over the half-line $\branch\ge0$ gives
$\hat\branch = \max(\tilde S_n,0)/\tilde\Fisher_n$ and
\[
\lambda_n \;=\; \big(\max(Z_n,0)\big)^2 + o_p(1).
\]
For $Z\sim\mathcal{N}(0,1)$:\ $\Prob(\max(Z,0)^2 \le x) = \Prob(Z\le 0) +
\Prob(0<Z\le\sqrt{x}) = \tfrac12 + \tfrac12\Prob(\chi^2_1\le x)$. The
mixture is exactly the geometry: with probability $\tfrac12$ the
unconstrained optimum wants negative excitation, the constrained optimum
sits at the boundary, and the LR is zero. When more than one parameter is
bounded (e.g.\ a nonnegative excitation matrix), the mixture weights
become the probabilities that the projected Gaussian lands on each face
of the cone \citep{selfliang1987}, and closed forms disappear ---
simulate instead.
\end{proof}

\emph{Practical rule.} The $p$-value for the boundary LRT is
$p = \tfrac12\,\Prob(\chi^2_1 \ge \lambda_{\mathrm{obs}})$ --- halve the
naive $p$-value; the $5\%$ critical value is $2.71$, not $3.84$. A team
using $\chi^2_1$ tables is running a conservative test and will
under-detect excitation; the direction of this error is at least safe,
which is why the mistake survives unnoticed.

\begin{warning}[When the decay is free, the mixture is wrong]\label{warn:ch07:boundary}
Under $H_0:\branch=0$ the decay $\decay$ is unidentified --- every
$\decay$ yields the same null model --- so if $\decay$ is estimated, the
LR is effectively a supremum over the decay bank and is stochastically
larger than \cref{eq:ch07:chibar}; the $50{:}50$ rule then
\emph{over}-detects excitation. Two legal responses: (i) test per fixed
decay-bank member (each obeys \cref{thm:ch07:chibar}) with a multiplicity
correction across the bank; (ii) parametric bootstrap --- simulate from
the fitted null NHPP, refit both models, and compare
$\lambda_{\mathrm{obs}}$ to the simulated null distribution, reusing the
envelope harness of \cref{prot:ch07:envelope}. Separately: Wald standard
errors from the inverse Hessian are invalid for any parameter whose
estimate sits at a boundary (fitted zeros in nonnegative excitation
matrices accumulate there); report profile-likelihood intervals or
held-out NLL instead.
\end{warning}

\section{Identifiability diagnostics}\label{ch07:sec:identifiability}

A model comparison is only as meaningful as the parameters it compares
are identified. Three diagnostics are standing obligations for the Hawkes
layer.

\paragraph{The $\branch$--$\decay$ ridge.}
Define the profile likelihood
$\loglik_p(\branch) = \max_{\decay,\psi}\loglik(\branch,\decay,\psi)$ and
its two-dimensional analogue over $(\branch,\decay)$. On event-time data
the profile is curved in both directions at sufficient volume; on
aggregated counts it is not, and the reason is elementary.

\begin{proposition}[Stationary counts pin $\branch$ and hide $\decay$]\label{prop:ch07:ridge}
For a stationary exponential Hawkes process with baseline $\mub$ and
kernel $\kernel(\tau)=\branch\decay e^{-\decay\tau}$, $\branch<1$, the
mean bucket count over width $\horizon$ is
$\Xi = \mub\,\horizon/(1-\branch)$. Hence
\[
\frac{\partial \Xi}{\partial\branch}
 = \frac{\mub\,\horizon}{(1-\branch)^2} > 0,
\qquad
\frac{\partial \Xi}{\partial\decay} = 0 :
\]
first moments of counts identify $\branch$ and carry \emph{zero}
information about $\decay$. All $\decay$-information in counts lives in
the temporal covariance structure, whose sensitivity to $\decay$
collapses at rate $\horizon^2 e^{-2\decay(1-\branch)\horizon}$ as the
bucket widens (\cref{ch:censoring}).
\end{proposition}

\begin{proof}
The stationary mean intensity solves the first-moment equation
$\bar\xi = \mub + \bar\xi\int_0^\infty \kernel(\tau)\dd\tau
 = \mub + \branch\,\bar\xi$ (\cref{ch:hawkes}), so
$\bar\xi = \mub/(1-\branch)$, and $\Xi = \bar\xi\horizon$ does not
involve $\decay$. Differentiation is immediate. The covariance statement
is \cref{ch:censoring}'s Fisher-information computation, restated here
because it is the operative rule: the exponential factor beats any
polynomial, so no amount of clever estimation recovers a timescale from
buckets much wider than the kernel's half-life.
\end{proof}

The consequences were measured (\cref{app:results}): $\decay$ estimated
from 28-day buckets came in biased low by ${\sim}20\%$ with a deceptively
tight spread --- the signature of an unidentified direction collapsing
onto a ridge --- while $\branch$ was stable at every width. The standing
rule, restated from \cref{ch:censoring}: bucket width at most one kernel
half-life is safe for timing; four or more half-lives destroys it; when
in the destroyed regime, fix $\decay$ to a decay bank and report
$\branch$ only. Operationally, every Hawkes fit ships its profile slices:
if $\loglik_p$ varies by less than $0.5$ nats across an order of
magnitude in $\decay$, the fit \emph{declares} $\decay$ unidentified and
the reporting layer refuses to print a decay estimate.

\paragraph{Frailty versus contagion: the upper-bound rule.}
Overdispersion and clustering are produced identically by contagion
(past deals raise future intensity) and by frailty (persistent latent
heterogeneity) \citep{duffie2009,azizpour2018}; a fitted Hawkes model
channels both into $\hat\branch$. The rehearsal quantified the
masquerade: under a latent AR($0.90$) frailty factor the fitted effective
excitation radius came out $2.4\times$ the truth (\cref{app:results}).
Since additional structure (shared covariates, factor models) can only
remove apparent excitation, every fitted branching ratio in this
repository is reported as an \emph{upper bound} on contagion --- the
sentence ``$\hat\branch=0.3$'' is banned; the sentence ``excitation plus
unmodeled common factors account for at most $\hat\branch=0.3$'' is the
template (\cref{ch:hawkes} develops the decomposition attempts and their
limits).

\paragraph{Parameter recovery as a shipping gate.}

\begin{protocol}[Parameter-recovery rehearsal]\label{prot:ch07:recovery}
No new model family ships until it passes, on synthetic data, all of:
\begin{enumerate}[label=(\alph*)]
\item \emph{Gradient check.} Analytic gradients of the likelihood agree
with finite differences at relative error $10^{-6}$--$10^{-8}$ at random
interior points --- the repository's standing discipline
(\cref{app:results}); a family whose gradients do not check does not
proceed to fitting.
\item \emph{Recovery.} Simulate $\ge 20$ datasets from the family at
parameter values matched to fitted real-data scale, pass them through
the synthetic observation layer with defects (C1)--(C6) active
(\cref{ch02:sec:synthetic}), fit with the production code path, and
verify: parameters claimed interpretable are recovered with bias small
against their sampling spread; parameters known unidentified under the
observation scheme (e.g.\ $\decay$ under coarse aggregation,
\cref{prop:ch07:ridge}) are \emph{declared} non-recoverable, and the
declaration is part of the family's documentation.
\item \emph{Coverage.} Whatever interval procedure the family advertises
(profile, bootstrap, posterior) achieves approximately nominal coverage
in the rehearsal.
\item \emph{Archival.} The rehearsal is an \code{experiments/} script
with logged seed and config, its results filed under the conventions of
\cref{prot:ch07:temporal}(d).
\end{enumerate}
A family failing (b) for some parameter may still ship \emph{for
prediction}, but that parameter is barred from interpretation until a
passing rehearsal exists. Defect exposure: the gate is exactly as honest
as the observation layer's reproduction of (C1)--(C6); (C3)'s unlabeled
exits are simulated, which is the one place we can rehearse a defect we
cannot observe.
\end{protocol}

\section{The temporal evaluation protocol}\label{ch07:sec:temporal}

All of the above presumes an evaluation era. Choosing it is not
bookkeeping: with 13.5 years of data spanning a zero-rate boom (2021), a
rate-shock bust (2022--2023), and a partial normalization (2024), the
split determines which regime the headline numbers describe.

\begin{protocol}[Temporal evaluation]\label{prot:ch07:temporal}
\leavevmode
\begin{enumerate}[label=(\alph*)]
\item \emph{Development split (fixed).} One split serves all model
development: train on data through 2019-12-31; validate (model selection,
hyperparameters, calibration maps) on 2020-01-01--2021-12-31; test on
2022-01-01--2024-06-30, touched only when a family is frozen. No result
computed on the test era before freezing may influence any modeling
choice.
\item \emph{Rolling-origin evaluation (headline).} Frozen families are
reported by rolling origin: expanding-window refits at quarterly steps
across the test era, each evaluated on the following quarter; the
headline number is the mean across origins with its range. Rolling
origins convert one number into a small distribution and are the only
defense against a verdict that is an artifact of a single split point.
\item \emph{No post-split information, anywhere.} Every transform ---
imputation, scaling, feature construction, probability recalibration ---
obeys \cref{prot:ch02:pit}: fitted on train (calibration maps on
validation), applied frozen downstream. Covariate support is checked:
test-era covariate ranges are reported next to every NLL, because a
log-linear intensity extrapolated outside the train support is
exponential extrapolation.
\item \emph{Reproducibility.} Every evaluation run logs its seed, config
hash, split boundaries, and code version; randomized-PIT draws and
bootstrap resamples are seeded; results are machine-readable files under
\code{results/}, one per experiment, following the \code{experiments/}
conventions.
\item \emph{Defect exposure.} (C1) is respected by construction (the
test era ends at the administrative censoring time and open spells carry
their survival factors); (C2)/(C3) exposure is inherited from the risk
sets used, and every reported metric names its risk-set rule
(\cref{ch07:sec:reporting}).
\end{enumerate}
\end{protocol}

\begin{decision}{D7.3 --- default temporal split}
The default development split is train $\le$ 2019, validation 2020--2021,
test 2022--June 2024, as in \cref{prot:ch07:temporal}(a). Rationale: the
validation era contains a violent regime change (COVID shock, then the
2021 boom), so model selection is stress-tested against nonstationarity
rather than sheltered from it; the test era contains the bust and
normalization, which is the regime a deployed model faces next; and
train ends before any of it, so no family is selected on the regime it
will be scored on. The known risk is that 2021--2023 spans a structural
break, so fixed-split verdicts may be regime artifacts --- which is
precisely why headlines come from rolling origins, not from the fixed
split. Reversal trigger: if rolling-origin family orderings disagree with
fixed-split validation orderings in two consecutive campaigns, the fixed
split is replaced by a multi-split (nested rolling) development scheme
and this box is rewritten.
\end{decision}

\section{Reporting discipline and defect exposure}\label{ch07:sec:reporting}

\begin{decision}{D7.4 --- no number without its caveat}
The following pairings are mandatory in every report, dashboard, README
table, and paper draft produced from this repository:
\begin{enumerate}[label=(\roman*)]
\item no hazard or intensity estimate without its \emph{risk-set rule}
(who was at risk, from when, closed by what);
\item no branching ratio without its \emph{upper-bound caveat}
(\cref{ch07:sec:identifiability});
\item no probability without its \emph{calibration plot} (reliability
diagram or PIT, per \cref{def:ch07:calibrated});
\item no ranking metric without its \emph{risk-set definition} (the
top-5 rate 0.14 vs 0.29 observed-vs-oracle gap of \cref{app:results} is
the standing proof that the risk set, not the model, can be the largest
term);
\item no count NLL without its \emph{dispersion} printed alongside
(\cref{prot:ch07:pit}).
\end{enumerate}
Rationale: each rule is the fossil of a specific rehearsal failure in
which the naked number was misleading and its caveat reversed the
conclusion. Reversal trigger: none anticipated; amendments add pairings,
they do not remove them.
\end{decision}

\paragraph{Defect declaration (per decision D2.1).}
The ``models'' of this chapter are the tests themselves, and they are
corrupted by the same six defects of \cref{tab:ch02:defects} they are
meant to police. (C1) \emph{robust}: administrative censoring is handled
exactly --- the open final rescaled interval is treated as censored
(\cref{prot:ch07:rescaling}b) and censored cells enter count likelihoods
with their survival factors. (C2) \emph{robust at sector level,
conditional at firm level}: all firm-level GOF statements are statements
about the entered (post-first-deal) universe; the protocols do not
correct the selection, they inherit \cref{ch:universe}'s scope. (C3)
\emph{vulnerable, optimistic} --- see \cref{warn:ch07:deadgof}. (C4)
\emph{corrects by resolution}: the batteries are run at the honest
resolution only --- weekly PIT always; event-time rescaling only on the
daily-resolved track --- so date jitter is absorbed rather than tested
against; the cost is that sub-weekly misfit is invisible by design. (C5)
\emph{vulnerable, detectable}: stale covariates push lack of fit into the
staleness-age strata of \cref{eq:ch07:mres} and into calendar PIT drift;
the diagnostics detect the damage but cannot attribute it without
\cref{ch:stale}'s machinery. (C6) \emph{robust for the ground process}
(sector counts are complete), \emph{thinned for the mark factor} (ranker
metrics run on tracked risk sets with the outside option of
\cref{ch:universe}).

\begin{warning}[GOF cannot detect dead exposure]\label{warn:ch07:deadgof}
Under (C3), dead firms held at risk deflate every fitted firm-level
hazard --- and the \emph{same} contaminated risk set is used at
evaluation time, so the deflated model is scored against the corruption
that produced it: held-out NLL, time-rescaling, and reliability can all
pass while the hazard is wrong for every living firm. Goodness of fit is
a consistency check between model and observed data, not between model
and truth. This is why risk-set truthfulness has its own diagnostics
(\cref{ch:poisson,ch:ranking}), why the exit-classifier AUC of $0.877$
was insufficient to repair it (\cref{app:results}), and why rule (i) of
decision D7.4 exists: the risk-set rule \emph{is} part of the estimate.
\end{warning}

\begin{codehooks}
Existing: \modrepo{mbpp/gof.py} implements
\code{time\_rescaling\_residuals}, \code{ks\_test\_exp1} (numpy-only KS),
\code{qq\_exp1}, \code{poisson\_pearson\_residuals}, and
\code{dispersion} --- the cores of
\cref{prot:ch07:rescaling,prot:ch07:pit}. To be written:
\modrepo{eval/}, the evaluation harness that makes the constitution
executable: \modrepo{eval/ladders.py} (held-out NLL versus the three
named baselines, nats per event and per cell, decision D7.2);
\modrepo{eval/pit.py} (randomized and non-randomized PIT, reliability
binning, PIT-drift series); \modrepo{eval/rescaling.py} (wraps
\code{mbpp.gof} with censored-final-gap handling, stratified martingale
residuals \cref{eq:ch07:mres}, bootstrap bands per
\cref{warn:ch07:estimated}); \modrepo{eval/envelope.py} (simulation
envelopes over \modrepo{eventtime/simulate.py} and
\modrepo{mbpp/ic\_simulate.py}, doubling as the bootstrap null for
\cref{warn:ch07:boundary}); \modrepo{eval/boundary.py} (chi-bar LRT,
\cref{thm:ch07:chibar}); \modrepo{eval/rolling.py} (the rolling-origin
driver of \cref{prot:ch07:temporal}). Conventions: every campaign is an
\code{experiments/expNN\_name.py} script emitting seeded, config-stamped
JSON to \code{results/}; the recovery gate of \cref{prot:ch07:recovery}
reuses \code{synthetic/} world generation; profile-likelihood slices and
the gradient check attach to the fitters in
\modrepo{sector\_stability.py}, \modrepo{eventtime/estimate.py}, and
\modrepo{mbpp/}.
\end{codehooks}
